
\documentclass{article}
\usepackage{colortbl}
\usepackage{makecell}
\usepackage{multirow}
\usepackage{supertabular}

\begin{document}

\newcounter{utterance}

\centering \large Interaction Transcript for game `referencegame', experiment `random\_grids', episode 11 with reflexion{-}Qwen3.5{-}4B.
\vspace{24pt}

{ \footnotesize  \setcounter{utterance}{1}
\setlength{\tabcolsep}{0pt}
\begin{supertabular}{c@{$\;$}|p{.15\linewidth}@{}p{.15\linewidth}p{.15\linewidth}p{.15\linewidth}p{.15\linewidth}p{.15\linewidth}}
    \# & $\;$A & \multicolumn{4}{c}{Game Master} & $\;\:$B\\
    \hline

    \theutterance \stepcounter{utterance}  
    & & \multicolumn{4}{p{0.6\linewidth}}{
        \cellcolor[rgb]{0.9,0.9,0.9}{
            \makecell[{{p{\linewidth}}}]{
                \texttt{\tiny{[P1$\langle$GM]}}
                \texttt{You are given three grids, where each of them is 5 by 5 in size.} \\
\texttt{Grids have empty cells marked with "▢" and filled cells marked with "X".} \\
\texttt{Your task is to generate a referring expression that best describes the target grid while distinguishing it from the two other distractor grids.} \\
\texttt{The first grid is the target grid, and the following two grids are the distractors.} \\
\\ 
\texttt{Target grid:} \\
\\ 
\texttt{X □ X X □} \\
\texttt{□ X □ X X} \\
\texttt{X X □ □ X} \\
\texttt{□ X X □ X} \\
\texttt{X □ X □ □} \\
\\ 
\texttt{Distractor grid 1:} \\
\\ 
\texttt{□ □ X X □} \\
\texttt{X □ □ X X} \\
\texttt{□ X X □ X} \\
\texttt{X X □ □ X} \\
\texttt{□ X X □ □} \\
\\ 
\texttt{Distractor grid 2:} \\
\\ 
\texttt{X X □ X □} \\
\texttt{□ X X □ X} \\
\texttt{X □ □ X □} \\
\texttt{□ X □ □ X} \\
\texttt{X □ X □ X} \\
\\ 
\texttt{Instruction: Describe the target grid.} \\
\texttt{Generate the referring expression starting with the tag "Expression: " for the given target grid. Omit any other text.} \\
            }
        }
    }
    & \\ \\

    \theutterance \stepcounter{utterance}  
    & \multicolumn{4}{p{0.6\linewidth}}{
        \cellcolor[rgb]{0.9,0.9,0.9}{
            \makecell[{{p{\linewidth}}}]{
                \texttt{\tiny{[P1$\rangle$GM]}}
                \texttt{Expression: The grid with an X in the top{-}left corner, an X in the top{-}right corner, and an X in the bottom{-}left corner, where the second row has Xs at the second and fourth columns, the third row has Xs at the first and fifth columns, the fourth row has Xs at the second and fifth columns, and the fifth row has Xs at the first and third columns.} \\
            }
        }
    }
    & & \\ \\

    \theutterance \stepcounter{utterance}  
    & & & \multicolumn{2}{p{0.3\linewidth}}{
        \cellcolor[rgb]{0.9,0.9,0.9}{
            \makecell[{{p{\linewidth}}}]{
                \texttt{\tiny{[GM$|$GM]}}
                \texttt{Expression: The grid with an X in the top{-}left corner, an X in the top{-}right corner, and an X in the bottom{-}left corner, where the second row has Xs at the second and fourth columns, the third row has Xs at the first and fifth columns, the fourth row has Xs at the second and fifth columns, and the fifth row has Xs at the first and third columns.} \\
            }
        }
    }
    & & \\ \\

    \theutterance \stepcounter{utterance}  
    & & \multicolumn{4}{p{0.6\linewidth}}{
        \cellcolor[rgb]{0.9,0.9,0.9}{
            \makecell[{{p{\linewidth}}}]{
                \texttt{\tiny{[GM$\rangle$P2]}}
                \texttt{You are given three grids, where each of them is 5 by 5 in size.} \\
\texttt{Grids have empty cells marked with "▢" and filled cells marked with "X".} \\
\texttt{You are also given a referring expression that describes one of the given grids.} \\
\texttt{Your task is to select the grid that matches the given referring expression.} \\
\texttt{Generate only the number (in text) of the grid that matches the given expression by selecting first, second, or third.} \\
\\ 
\texttt{First:} \\
\\ 
\texttt{X X □ X □} \\
\texttt{□ X X □ X} \\
\texttt{X □ □ X □} \\
\texttt{□ X □ □ X} \\
\texttt{X □ X □ X} \\
\\ 
\texttt{Second:} \\
\\ 
\texttt{□ □ X X □} \\
\texttt{X □ □ X X} \\
\texttt{□ X X □ X} \\
\texttt{X X □ □ X} \\
\texttt{□ X X □ □} \\
\\ 
\texttt{Third:} \\
\\ 
\texttt{X □ X X □} \\
\texttt{□ X □ X X} \\
\texttt{X X □ □ X} \\
\texttt{□ X X □ X} \\
\texttt{X □ X □ □} \\
\\ 
\texttt{Expression: The grid with an X in the top{-}left corner, an X in the top{-}right corner, and an X in the bottom{-}left corner, where the second row has Xs at the second and fourth columns, the third row has Xs at the first and fifth columns, the fourth row has Xs at the second and fifth columns, and the fifth row has Xs at the first and third columns.} \\
\texttt{Question: Which grid does the expression refer to?} \\
\texttt{Start with the tag "Answer: ", followed by your selection. Omit any other text.} \\
            }
        }
    }
    & \\ \\

    \theutterance \stepcounter{utterance}  
    & & & \multicolumn{4}{p{0.6\linewidth}}{
        \cellcolor[rgb]{0.9,0.9,0.9}{
            \makecell[{{p{\linewidth}}}]{
                \texttt{\tiny{[GM$\langle$P2]}}
                \texttt{Answer: First} \\
            }
        }
    }
     \\ \\

    \theutterance \stepcounter{utterance}  
    & & & \multicolumn{2}{p{0.3\linewidth}}{
        \cellcolor[rgb]{0.9,0.9,0.9}{
            \makecell[{{p{\linewidth}}}]{
                \texttt{\tiny{[GM$|$GM]}}
                \texttt{first} \\
            }
        }
    }
    & & \\ \\

    \theutterance \stepcounter{utterance}  
    & \multicolumn{4}{p{0.6\linewidth}}{
        \cellcolor[rgb]{0.95,0.95,0.95}{
            \makecell[{{p{\linewidth}}}]{
                \texttt{\tiny{[P2$\rightarrow$P1]}}
                \texttt{Post{-}episode Reflexion (not sent to GM):} \\
\texttt{New plan: I failed because I described the target grid using a generic, exhaustive list of coordinates without explicitly referencing the unique features that distinguish it from the distractors. Specifically, I missed the most critical differentiator: the target grid is the **only** one with an 'X' in the top{-}left corner (Distractor 1 has a '▢' there, and Distractor 2 has an 'X' but lacks the specific combination in the top row). I should have started the expression by highlighting this unique top{-}left 'X' to immediately separate the target from Distractor 1, and then noted the specific pattern in the third row (Xs at columns 1 and 5) which distinguishes it from Distractor 2 (which has Xs at columns 1 and 4 in that row). The new strategy is to identify and lead with the single most distinctive feature of the target grid relative to the distractors, rather than listing all cells.} \\
            }
        }
    }
    & & \\ \\

    \theutterance \stepcounter{utterance}  
    & & & \multicolumn{4}{p{0.6\linewidth}}{
        \cellcolor[rgb]{0.95,0.95,0.95}{
            \makecell[{{p{\linewidth}}}]{
                \texttt{\tiny{[P1$\rightarrow$P2]}}
                \texttt{Post{-}episode Reflexion (not sent to GM):} \\
\texttt{New plan: I incorrectly identified the target grid by relying on a partial match of the description rather than verifying every specific coordinate constraint. The expression requires an X in the top{-}right corner and specific patterns in rows two through five, which the "First" grid does not satisfy (e.g., it lacks an X in the top{-}right). Next time, I will systematically check each grid against the full list of constraints: first verifying the top{-}left and top{-}right corners, then confirming the exact column positions of Xs in rows 2, 3, 4, and 5. Only if all conditions match perfectly will I select that grid number.} \\
            }
        }
    }
     \\ \\

\end{supertabular}
}

\end{document}
